\subsection{Mockups}\label{sec:mockups}

Auf Basis der in \autoref{sec:skizzen_wireframes} validierten Wireframes wurden die finalen Mockups für beide Varianten (\autoref{subsec:varianten}) entwickelt.

\subsubsection{Variante~1: Desktop-Variante}

Die Desktop-Variante ist für Querformat ausgelegt und wird auf zwei Bildschirmgrößen präsentiert: 1080p-Desktop und iPad-Querformat. Sie nutzt das in \autoref{sec:skizzen_wireframes} entwickelte Vier-Spalten-Layout (Drawer, SideNav, Master, Detail) zur optimalen Nutzung der horizontalen Bildschirmfläche.

\subsubsection*{Datei-Explorer}

\paragraph*{Ist-Zustand (Vorher)}
Im Ist-Zustand (\autoref{fig:ist_dateiexplorer_ohne_hierarchie}) war die Ordnerhierarchie in einer flachen Liste dargestellt \textendash{} im Usability-Test verloren Nutzer:innen den Überblick. Die Dateiauswahl (\autoref{fig:ist_dateiexplorer_dateiauswahl}) war mit einem unstrukturierten Detailbereich versehen.

\begin{figure}[H]
    \centering
    \includegraphics[width=0.9\textwidth]{images/3.3_Mockups/Ist-Desktop/Ist-Dateiexplorer-Ordnerauswahl.png}
    \caption[Ist-Zustand: Datei-Explorer (Vorher)]{Ist-Zustand -- Datei-Explorer (Ordnerauswahl). Quelle: Eigene Darstellung.}
    \label{fig:ist_dateiexplorer_ohne_hierarchie}
\end{figure}

\begin{figure}[H]
    \centering
    \includegraphics[width=0.9\textwidth]{images/3.3_Mockups/Ist-Desktop/Ist-Dateiexplorer-Dateiauswahl.png}
    \caption[Ist-Zustand: Datei-Explorer (Dateiauswahl)]{Ist-Zustand -- Datei-Explorer (Dateiauswahl mit ausgeklappten Details). Quelle: Eigene Darstellung.}
    \label{fig:ist_dateiexplorer_dateiauswahl}
\end{figure}

\paragraph*{Redesign (Nachher)}
\autoref{fig:mockup_desktop_dateiexplorer_ordnerauswahl_desktop} zeigt die Ordnerauswahl mit Detail-Ansicht auf Desktop 1080p, \autoref{fig:mockup_desktop_dateiexplorer_ordnerauswahl_ipad_quer} im iPad-Querformat.

\begin{figure}[H]
    \centering
    \includegraphics[width=0.9\textwidth]{images/3.3_Mockups/Desktop-1080p/Datei-Explorer-Ordnerauswahl-Stammverzeichnis.png}
    \caption[Desktop-Variante: Datei-Explorer (Ordnerauswahl, Desktop 1080p)]{Desktop-Variante -- Datei-Explorer (Ordnerauswahl, Desktop 1080p). Quelle: Eigene Darstellung.}
    \label{fig:mockup_desktop_dateiexplorer_ordnerauswahl_desktop}
\end{figure}

\begin{figure}[H]
    \centering
    \includegraphics[width=0.9\textwidth]{images/3.3_Mockups/Desktop-iPad-Quer/Datei-Explorer-Ordnerauswahl-Stammverzeichnis.png}
    \caption[Desktop-Variante: Datei-Explorer (Ordnerauswahl, iPad Quer)]{Desktop-Variante -- Datei-Explorer (Ordnerauswahl, iPad Querformat). Quelle: Eigene Darstellung.}
    \label{fig:mockup_desktop_dateiexplorer_ordnerauswahl_ipad_quer}
\end{figure}

\newpage

\autoref{fig:mockup_desktop_dateiexplorer_dateiauswahl} zeigt die Dateiauswahl mit Detail-Ansicht auf Desktop 1080p, \autoref{fig:mockup_desktop_dateiexplorer_dateiauswahl_ipad} im iPad-Querformat. Die Breadcrumb-Navigation verdeutlicht die Position des Nutzers in der Ordnerhierarchie.

\begin{figure}[H]
    \centering
    \includegraphics[width=0.9\textwidth]{images/3.3_Mockups/Desktop-1080p/Datei-Explorer-Dateiauswahl.png}
    \caption[Desktop-Variante: Datei-Explorer (Dateiauswahl, Desktop 1080p)]{Desktop-Variante -- Datei-Explorer (Dateiauswahl, Desktop 1080p). Quelle: Eigene Darstellung.}
    \label{fig:mockup_desktop_dateiexplorer_dateiauswahl}
\end{figure}

\begin{figure}[H]
    \centering
    \includegraphics[width=0.9\textwidth]{images/3.3_Mockups/Desktop-iPad-Quer/Datei-Explorer-Dateiauswahl.png}
    \caption[Desktop-Variante: Datei-Explorer (Dateiauswahl, iPad Quer)]{Desktop-Variante -- Datei-Explorer (Dateiauswahl, iPad Querformat). Quelle: Eigene Darstellung.}
    \label{fig:mockup_desktop_dateiexplorer_dateiauswahl_ipad}
\end{figure}

\newpage

\subsubsection*{Öffentliche Ordneransicht}

\paragraph*{Ist-Zustand (Vorher)}
Im Ist-Zustand (\autoref{fig:ist_anonymous_dateiauswahl}) wurde bei Dateiauswahl in der öffentlichen Ordneransicht der Download-Button direkt in der Dateiliste angezeigt \textendash{} zusätzlich zur nicht deutlich erkennbaren Sortierung, die im Usability-Test zu Unsicherheit führte.

\begin{figure}[H]
    \centering
    \includegraphics[width=0.9\textwidth]{images/3.3_Mockups/Ist-Desktop/Ist-Oeffentliche-Ordneransicht-Dateiauswahl.png}
    \caption[Ist-Zustand: Öffentliche Ordnerfreigabe (Vorher)]{Ist-Zustand -- Öffentliche Ordnerfreigabe. Quelle: Eigene Darstellung.}
    \label{fig:ist_anonymous_dateiauswahl}
\end{figure}

\vfill

\paragraph*{Redesign (Nachher)}

\autoref{fig:mockup_desktop_anonymous_dateiauswahl_desktop} zeigt die öffentliche Ordneransicht mit Dateiauswahl auf Desktop 1080p, \autoref{fig:mockup_desktop_anonymous_dateiauswahl_ipad_quer} im iPad-Querformat ohne Dateiauswahl.

\begin{figure}[H]
    \centering
    \includegraphics[width=0.9\textwidth]{images/3.3_Mockups/Desktop-1080p/Oeffentliche-Ordnerfreigabe-Dateiauswahl.png}
    \caption[Desktop-Variante: Öffentliche Ordneransicht (Dateiauswahl, Desktop 1080p)]{Desktop-Variante -- Öffentliche Ordnerfreigabe (Dateiauswahl, Desktop 1080p). Quelle: Eigene Darstellung.}
    \label{fig:mockup_desktop_anonymous_dateiauswahl_desktop}
\end{figure}

\begin{figure}[H]
    \centering
    \includegraphics[width=0.9\textwidth]{images/3.3_Mockups/Desktop-iPad-Quer/Oeffentliche-Ordnerfreigabe-Ohne-Auswahl.png}
    \caption[Desktop-Variante: Öffentliche Ordneransicht (Ohne Auswahl, iPad Quer)]{Desktop-Variante -- Öffentliche Ordnerfreigabe (Ohne Auswahl, iPad Querformat). Quelle: Eigene Darstellung.}
    \label{fig:mockup_desktop_anonymous_dateiauswahl_ipad_quer}
\end{figure}

\subsubsection*{Einstellungen}

\paragraph*{Ist-Zustand (Vorher)}
In den Einstellungen fehlte im Ist-Zustand der fixierte Footer. Die Speichern-Schaltfläche war somit vom Content-Overflow betroffen.

\begin{figure}[H]
    \centering
    \includegraphics[width=0.9\textwidth]{images/3.3_Mockups/Ist-Desktop/Ist-Einstellungen-Clean.png}
    \caption[Ist-Zustand: Einstellungen (Vorher)]{Ist-Zustand (keine Footer-Sektion). Quelle: Eigene Darstellung.}
    \label{fig:ist_einstellungen_cleans}
\end{figure}

\paragraph*{Redesign (Nachher)}\null\hfill
Das neue Einstellungsmenü besitzt nun ein Clean-Dirty-State-Management. Der Footer ist fixiert und enthält die Speichern-Schaltfläche die bei einer Änderung (Dirty) aktiv wird.

\begin{figure}[H]
    \centering
    \includegraphics[width=0.9\textwidth]{images/3.3_Mockups/Desktop-1080p/Einstellungen-Clean.png}
    \caption[Desktop-Variante: Einstellungen (Clean-Zustand, Desktop 1080p)]{Desktop-Variante -- Einstellungen (Clean-Zustand, Desktop 1080p). Quelle: Eigene Darstellung.}
    \label{fig:mockup_desktop_einstellungen_clean_desktop}
\end{figure}

\begin{figure}[H]
    \centering
    \includegraphics[width=0.9\textwidth]{images/3.3_Mockups/Desktop-iPad-Quer/Einstellungen.png}
    \caption[Desktop-Variante: Einstellungen (Clean-Zustand, iPad Quer)]{Desktop-Variante -- Einstellungen (Clean-Zustand, iPad Querformat). Quelle: Eigene Darstellung.}
    \label{fig:mockup_desktop_einstellungen_clean_ipad_quer}
\end{figure}

\begin{figure}[H]
    \centering
    \includegraphics[width=0.9\textwidth]{images/3.3_Mockups/Desktop-1080p/Einstellungen-Dirty.png}
    \caption[Desktop-Variante: Einstellungen (Dirty-State)]{Desktop-Variante -- Einstellungen (Dirty-State mit hervorgehobenem Speichern-Button). Quelle: Eigene Darstellung.}
    \label{fig:mockup_desktop_einstellungen_dirty}
\end{figure}

\begin{figure}[H]
    \centering
    \includegraphics[width=0.9\textwidth]{images/3.3_Mockups/Desktop-1080p/Einstellungen-Dirty-Navigation.png}
    \caption[Desktop-Variante: Einstellungen (Navigations-Schutzabfrage)]{Desktop-Variante -- Einstellungen (Navigations-Schutzabfrage). Quelle: Eigene Darstellung.}
    \label{fig:mockup_desktop_einstellungen_dirty_navigation}
\end{figure}

\begin{figure}[H]
    \centering
    \includegraphics[width=0.9\textwidth]{images/3.3_Mockups/Desktop-iPad-Quer/Einstellungen-Collapsed-Drawer.png}
    \caption[Desktop-Variante: Einstellungen (iPad Collapsed)]{Desktop-Variante -- Einstellungen (iPad Querformat, Collapsed Drawer). Quelle: Eigene Darstellung.}
    \label{fig:mockup_desktop_einstellungen_ipad_quer}
\end{figure}

\newpage

\subsubsection{Variante~2: Mobil-Variante}\label{subsec:Mobil_Variante}

Die Mobil-Variante ersetzt das Desktop-Layout durch ein ein-spaltiges Design mit Hamburger-Menü und horizontaler Filterleiste (siehe \autoref{sec:skizzen_wireframes} für die Layout-Begründung). Touch-Targets sind für die Finger-Bedienung vergrößert. Die Variante ist für schmale Bildschirmformate (Hochformat) ausgelegt und wird auf zwei Bildschirmgrößen präsentiert: iPad-Hochformat und iPhone. Die konzeptionellen Navigations- und Layout-Unterschiede gegenüber der Desktop-Variante sind in \autoref{sec:informationsarchitektur} und \autoref{sec:skizzen_wireframes} dokumentiert. Die Detail-Sektion innerhalb des Datei-Explorers öffnet sich hier als Overlay und nimmt Vollbild ein.

\subsubsection*{Datei-Explorer}
\paragraph*{Ist-Zustand (Vorher)}

\autoref{fig:ist_mobil_dateiexplorer_drawer} zeigt den Ist-Zustand der Mobil-Variante. Die Drawer-Navigation ist teils unklar, die Touch-Targets zu klein, die Seite ist mit zu vielen Schaltflächen und Elementen versehen.

\begin{figure}[H]
    \centering
    \includegraphics[height=0.42\textheight, keepaspectratio]{images/3.3_Mockups/Ist-Mobil/Ist-Dateiexplorer-Drawer.png}\hspace{2em}
    \includegraphics[height=0.42\textheight, keepaspectratio]{images/3.3_Mockups/Ist-Mobil/Ist-Dateiexplorer-Dateiauswahl.png}
    \caption[Ist-Zustand: Datei-Explorer Mobil (Vorher)]{Ist-Zustand Mobil (links: Drawer mit unklarer Navigation, rechts: zu viele Schaltflächen, Touch-Targets zu klein). Quelle: Eigene Darstellung.}
    \label{fig:ist_mobil_dateiexplorer_drawer}
\end{figure}

\vfill

\paragraph*{Redesign (Nachher)}
\autoref{fig:mockup_mobil_dateiexplorer_ordnerauswahl} zeigt die Ordnerauswahl mit Detail-Ansicht auf iPhone und iPad Hochformat. Der Drawer ist klar strukturiert (\autoref{fig:mockup_mobil_dateiexplorer_expanded_drawer}), die Sekundärnavigation als horizontale Filterleiste umgesetzt. Tabellenzeilen und Schaltflächen sind vergrößert.

\begin{figure}[H]
    \centering
    \includegraphics[height=0.42\textheight, keepaspectratio]{images/3.3_Mockups/Mobil-iPhone/Datei-Explorer-Ordnernavigation.png}
    \includegraphics[height=0.42\textheight, keepaspectratio]{images/3.3_Mockups/Mobil-iPad-Hoch/Datei-Explorer-Ordnernavigation.png}\hspace{2em}
    \caption[Mobil-Variante: Datei-Explorer (Ordnerauswahl + Navigation)]{Mobil-Variante -- Datei-Explorer Ordnerauswahl (links: iPhone, rechts: iPad Hochformat). Quelle: Eigene Darstellung.}
    \label{fig:mockup_mobil_dateiexplorer_ordnerauswahl}
\end{figure}

\begin{figure}[H]
    \centering
    \includegraphics[height=0.38\textheight, keepaspectratio]{images/3.3_Mockups/Mobil-iPhone/Datei-Explorer-Expanded-Drawer.png}\hspace{2em}
    \caption[Mobil-Variante: Datei-Explorer (Expanded Drawer, iPhone)]{Mobil-Variante -- Datei-Explorer (Expanded Drawer, iPhone). Quelle: Eigene Darstellung.}
    \label{fig:mockup_mobil_dateiexplorer_expanded_drawer}
\end{figure}

\newpage
\autoref{fig:mockup_mobil_dateiexplorer_dateiauswahl} zeigt die Dateiauswahl mit Detail-Ansicht auf iPhone. Diese öffnet sich als Overlay und nimmt Vollbild ein.

\begin{figure}[H]
    \centering
    \includegraphics[height=0.42\textheight, keepaspectratio]{images/3.3_Mockups/Mobil-iPhone/Datei-Explorer-Dateiauswahl.png}\hspace{2em}
    \caption[Mobil-Variante: Datei-Explorer (Dateiauswahl)]{Mobil-Variante -- Datei-Explorer Dateiauswahl (iPhone). Quelle: Eigene Darstellung.}
    \label{fig:mockup_mobil_dateiexplorer_dateiauswahl}
\end{figure}

\vfill

\subsubsection*{Öffentliche Ordneransicht}

\paragraph*{Ist-Zustand (Vorher)}
Im Ist-Zustand erschienen Download-Buttons direkt in der Dateiliste ohne erkenntlich sortierter Liste, was im Usability-Test zu Unsicherheit führte. Ein Header mit Logo und Schließen-Button ermöglichte das Verlassen bei Dateiauswahl.

\begin{figure}[H]
    \centering
    \includegraphics[height=0.42\textheight, keepaspectratio]{images/3.3_Mockups/Ist-Mobil/Ist-Oeffentliche-Ordneransicht.png}\hspace{2em}
    \includegraphics[height=0.42\textheight, keepaspectratio]{images/3.3_Mockups/Ist-Mobil/Ist-Oeffentliche-Ordneransicht-Dateiauswahl.png}
    \caption[Ist-Zustand: Öffentliche Ordneransicht Mobil (Vorher)]{Ist-Zustand Mobil Öffentliche Ordneransicht (links: Dateiliste ohne Auswahl, rechts: Download-Button in der Dateiliste). Quelle: Eigene Darstellung.}
    \label{fig:ist_mobil_anonymous}
\end{figure}

\vfill

\paragraph*{Redesign (Nachher)}
\autoref{fig:mockup_mobil_ohne_auswahl_dateiauswahl} zeigt die öffentliche Ordneransicht mit aktiver Sortierung und ohne Dateiauswahl auf iPhone und iPad Hochformat. Die Dateiauswahl öffnet sich als Overlay und nimmt Vollbild ein (\autoref{fig:mockup_mobil_anonymous_dateiauswahl}).

\begin{figure}[H]
    \centering
    \includegraphics[height=0.42\textheight, keepaspectratio]{images/3.3_Mockups/Mobil-iPad-Hoch/Oeffentliche-Ordnerfreigabe-Ohne-Auswahl.png}\hspace{2em}
    \includegraphics[height=0.42\textheight, keepaspectratio]{images/3.3_Mockups/Mobil-iPhone/Oeffentliche-Ordnerfreigabe-Ohne-Auswahl.png}
    \caption[Mobil-Variante: Öffentliche Ordneransicht (Ohne Auswahl)]{Mobil-Variante -- Öffentliche Ordnerfreigabe Ohne Auswahl (links: iPad, rechts: iPhone). Quelle: Eigene Darstellung.}
    \label{fig:mockup_mobil_ohne_auswahl_dateiauswahl}
\end{figure}

\begin{figure}[H]
    \centering
    \includegraphics[height=0.40\textheight, keepaspectratio]{images/3.3_Mockups/Mobil-iPhone/Oeffentliche-Ordnerfreigabe-Dateiauswahl.png}
    \caption[Mobil-Variante: Öffentliche Ordnerfreigabe (Dateiauswahl iPhone)]{Mobil-Variante -- Datei-Explorer (Dateiauswahl, iPhone). Quelle: Eigene Darstellung.}
    \label{fig:mockup_mobil_anonymous_dateiauswahl}
\end{figure}

\subsubsection*{Einstellungen}

\paragraph*{Ist-Zustand (Vorher)}
In den Einstellungen fehlte die klare Trennung von Content und Footer.

\begin{figure}[H]
    \centering
    \includegraphics[height=0.42\textheight, keepaspectratio]{images/3.3_Mockups/Ist-Mobil/Ist-Einstellungen-Clean.png}\hspace{2em}
    \includegraphics[height=0.42\textheight, keepaspectratio]{images/3.3_Mockups/Ist-Mobil/Ist-Einstellungen-Drawer.png}
    \caption[Ist-Zustand: Einstellungen Mobil (Vorher)]{Ist-Zustand Mobil Einstellungen (links: keine Header/Footer-Trennung, rechts: inkonsistenter Drawer). Quelle: Eigene Darstellung.}
    \label{fig:ist_mobil_einstellungen}
\end{figure}

\vfill

\paragraph*{Redesign (Nachher)}
\autoref{fig:mockup_mobil_einstellungen_ipad_hoch} zeigt die Einstellungen mit fixiertem Footer auf iPad Hochformat, \autoref{fig:mockup_mobil_einstellungen_iphone} im iPhone-Format. Der Sekundärnavigation integriert sich in den Drawer. Die Speichern-Schaltfläche wird bei Änderungen (Dirty) aktiv.

\begin{figure}[H]
    \centering
    \includegraphics[height=0.40\textheight, keepaspectratio]{images/3.3_Mockups/Mobil-iPad-Hoch/Einstellungen.png}\hspace{2em}
    \includegraphics[height=0.40\textheight, keepaspectratio]{images/3.3_Mockups/Mobil-iPad-Hoch/Einstellungen-Collapsed-Drawer.png}
    \caption[Mobil-Variante: Einstellungen (iPad Hochformat)]{Mobil-Variante -- Einstellungen (links: iPad Hochformat, rechts: Collapsed Drawer). Quelle: Eigene Darstellung.}
    \label{fig:mockup_mobil_einstellungen_ipad_hoch}
\end{figure}

\begin{figure}[H]
    \centering
    \includegraphics[height=0.40\textheight, keepaspectratio]{images/3.3_Mockups/Mobil-iPhone/Einstellungen-Expanded-Drawer.png}\hspace{2em}
    \includegraphics[height=0.40\textheight, keepaspectratio]{images/3.3_Mockups/Mobil-iPhone/Einstellungen.png}
    \caption[Mobil-Variante: Einstellungen (iPhone)]{Mobil-Variante -- Einstellungen (links: iPhone Expanded Drawer, rechts: Drawer geschlossen). Quelle: Eigene Darstellung.}
    \label{fig:mockup_mobil_einstellungen_iphone}
\end{figure}
